%==============================================================================
%  GENERIC entálpico PARA DUMMIES
%  Explicación ab-initio, ecuación a ecuación, del modelo de soldadura,
%  el modelo FEM, la entalpía, las c_p, y la formulación GENERIC en entalpía.
%
%  Requisito de partida del lector: solo el 1er y 2do principio de la
%  termodinámica. NO se asume soldadura, ecuación del calor, FEM ni termo
%  avanzada.
%==============================================================================
\documentclass[11pt]{article}

\usepackage[utf8]{inputenc}
\usepackage[T1]{fontenc}
\usepackage[spanish,es-noshorthands]{babel}
\usepackage[margin=1in]{geometry}
\usepackage{amsmath,amssymb,amsthm,mathtools}
\usepackage{bm}
\usepackage{siunitx}
\sisetup{output-decimal-marker={.}}
\usepackage{graphicx}
\usepackage{booktabs}
\usepackage{array}
\usepackage{xcolor}
\usepackage{enumitem}
\usepackage[hidelinks]{hyperref}

% --- macros -------------------------------------------------------------------
\newcommand{\dd}{\mathrm{d}}
\newcommand{\grad}{\nabla}
\newcommand{\Tinf}{T_{\infty}}
\newcommand{\cpapp}{c_p^{\mathrm{app}}}
\newcommand{\cpsens}{c_p^{\mathrm{sens}}}
\newcommand{\Lw}{\bm{L}_{w}}
\newcommand{\bmu}{\bm{\mu}}
\newcommand{\abs}[1]{\left\lvert#1\right\rvert}
\newcommand{\norm}[1]{\left\lVert#1\right\rVert}
\newcommand{\alert}[1]{\textcolor{red}{\textbf{#1}}}

% --- cajas pedagógicas (entornos tipo teorema, siempre disponibles) -----------
\theoremstyle{definition}
\newtheorem{definicion}{Definición}[section]
\newtheorem{idea}{Idea clave}[section]
\newtheorem{ejemplo}{Ejemplo}[section]
\theoremstyle{plain}
\newtheorem{proposicion}{Proposición}[section]

% caja de aviso "para dummies"
\newcommand{\dummies}[1]{%
  \par\smallskip\noindent
  \colorbox{yellow!18}{\parbox{\dimexpr\linewidth-2\fboxsep}{%
  \small\textbf{En cristiano:} #1}}\par\smallskip}

\title{\bfseries GENERIC entálpico \emph{para dummies}\\[4pt]
  \large De la soldadura y el calor a una red neuronal que respeta\\
  el primer y el segundo principio, ecuación a ecuación}
\author{}
\date{}

\begin{document}
\maketitle

\begin{abstract}
\noindent
Este documento explica, \textbf{desde cero}, la formulación matemática que hay
detrás del modelo de soldadura del proyecto. La única herramienta previa que se
supone es el \textbf{primer y el segundo principio de la termodinámica}. A partir
de ahí construimos, paso a paso y ecuación a ecuación: (1)~qué pasa físicamente al
soldar; (2)~qué es el calor específico $c_p$ y por qué la \emph{entalpía} es la
variable correcta; (3)~la ecuación del calor, deducida de un balance de energía;
(4)~qué es el método de los elementos finitos (FEM) y por qué lo necesitamos;
(5)~qué es GENERIC y cómo garantiza los dos principios \emph{por construcción};
(6)~por qué imponer GENERIC sobre la \emph{temperatura} está mal cuando hay cambio
de fase, y (7)~cómo se arregla pasándolo a \emph{entalpía}. Cada símbolo se define
la primera vez que aparece. Las cajas amarillas resumen ``en cristiano'' lo que
cada bloque significa.
\end{abstract}

\tableofcontents
\bigskip

%==============================================================================
\section{Antes de empezar: lo único que damos por sabido}
%==============================================================================

Vamos a apoyarnos solo en dos frases. Todo lo demás se construye encima.

\begin{idea}[Primer principio: la energía ni se crea ni se destruye]
La energía de un sistema solo cambia si \emph{entra} o \emph{sale} energía a través
de su frontera. No aparece de la nada ni desaparece. Si nadie aporta ni quita
energía, la energía total es constante.
\end{idea}

\begin{idea}[Segundo principio: el desorden (entropía) no disminuye solo]
Existe una magnitud llamada \emph{entropía} $S$ que mide, en cierto sentido, el
``desorden'' o la irreversibilidad. En un sistema aislado la entropía
\textbf{nunca decrece}: $\dd S/\dd t \ge 0$. Una consecuencia cotidiana: el calor
fluye \textbf{espontáneamente de lo caliente a lo frío}, nunca al revés por sí
solo.
\end{idea}

Necesitaremos además una relación cuantitativa entre energía, calor y entropía que
se deduce de estos principios y que es la \emph{bisagra} de todo el documento. La
introducimos ya, y la justificamos en la Sección~\ref{sec:bisagra}:
\begin{equation}
  \boxed{\;\dd S = \frac{\dd E}{T}\;}
  \qquad\Longleftrightarrow\qquad
  \frac{\partial S}{\partial E} = \frac{1}{T}.
  \label{eq:bisagra}
\end{equation}
Aquí $E$ es la energía (en julios, $\mathrm{J}$), $T$ la temperatura absoluta (en
kelvin, $\mathrm{K}$) y $S$ la entropía ($\mathrm{J/K}$).

\dummies{Si a un cuerpo le metes una pizca de energía $\dd E$ estando a temperatura
$T$, su entropía sube en $\dd E/T$. Cuanto más caliente está, menos entropía gana
por el mismo julio. Esta fórmula es \emph{toda} la termodinámica que usaremos.}

%==============================================================================
\section{El problema físico: ¿qué pasa cuando soldamos?}
\label{sec:soldadura}
%==============================================================================

Imagina una \textbf{placa metálica} plana (acero), como una baldosa de metal.
Por encima pasa un \textbf{arco eléctrico} (la ``antorcha'' de soldadura) que se
mueve siguiendo un camino. El arco es una fuente de calor muy concentrada y muy
intensa: deposita energía en una zona pequeña que va recorriendo la placa.

Lo que ocurre, en orden:
\begin{enumerate}[leftmargin=1.5em]
  \item Justo bajo el arco el metal se calienta muchísimo, se \textbf{funde}
  (forma un ``charco'' o \emph{baño} de fusión) e incluso puede empezar a
  \textbf{hervir}/vaporizarse en el punto más caliente.
  \item El calor se \textbf{difunde} hacia el metal frío de alrededor: la zona
  caliente se ensancha como una mancha que se va suavizando.
  \item Cuando el arco se aleja, esa zona se \textbf{enfría}: cede calor al resto
  de la placa y al aire (convección y radiación).
\end{enumerate}

Lo que queremos predecir es el \textbf{campo de temperatura} $T(\bm x,t)$: la
temperatura en cada punto $\bm x$ de la placa y en cada instante $t$. Saber esto es
clave en ingeniería porque de ahí salen las tensiones residuales, las deformaciones
y la microestructura del material.

\dummies{Soldar $=$ pasear una fuente de calor brutal por una placa. El metal se
funde donde pasa y se enfría detrás. Queremos un mapa de temperaturas en cada
instante.}

\paragraph{¿Por qué es difícil?} Por dos motivos que dan nombre a casi todo lo que
sigue:
\begin{itemize}[leftmargin=1.5em]
  \item \textbf{Cambio de fase.} Fundir y hervir \emph{absorben} mucha energía
  \textbf{sin que suba la temperatura} (luego veremos por qué). El metal no se
  comporta como ``más energía $=$ más temperatura''.
  \item \textbf{Hay que avanzar en el tiempo muchísimos pasos.} La solución se
  calcula instante a instante; los errores se pueden ir acumulando.
\end{itemize}

%==============================================================================
\section{Energía, calor específico y por qué nace la entalpía}
\label{sec:cp}
%==============================================================================

\subsection{Calor específico: cuánta energía cuesta calentar}

\begin{definicion}[Calor específico $c_p$]
El \emph{calor específico} $c_p$ es la energía que hay que dar a $1\,\mathrm{kg}$ de
material para subir su temperatura $1\,\mathrm{K}$. Unidades:
$\mathrm{J\,kg^{-1}\,K^{-1}}$.
\end{definicion}

Para el acero, lejos de cambios de fase, $c_p \approx \SI{600}{J\,kg^{-1}\,K^{-1}}$.
Es decir, calentar un kilo de acero un grado cuesta unos 600 julios.

Si tengo una masa $m$ y le subo la temperatura de $T_0$ a $T$, la energía que he
metido (suponiendo $c_p$ constante) es
\begin{equation}
  E = m\,c_p\,(T-T_0).
  \label{eq:Esimple}
\end{equation}
Si $c_p$ \emph{no} es constante (depende de $T$), hay que sumar pizca a pizca, es
decir integrar:
\begin{equation}
  E = m\int_{T_0}^{T} c_p(\tau)\,\dd\tau .
  \label{eq:Eint}
\end{equation}

\dummies{$c_p$ es el ``precio en julios'' de subir un grado un kilo. La energía
almacenada es ese precio multiplicado por los grados que has subido (y si el precio
cambia con la temperatura, integras).}

Como trabajamos con campos (energía \emph{por unidad de volumen}) en vez de masas
sueltas, usamos la \textbf{densidad} $\rho$ (en $\mathrm{kg/m^3}$; para el acero
$\rho=\SI{7850}{kg/m^3}$). La masa de un trocito de volumen $V$ es $m=\rho V$, así
que la energía por unidad de volumen es $E/V = \rho\int c_p\,\dd T$.

\subsection{El calor latente: energía que entra y la temperatura no se mueve}

Aquí está la rareza de la soldadura. Cuando el metal \textbf{funde}, los átomos
rompen los enlaces del sólido para volverse líquido. Eso \emph{cuesta energía},
pero esa energía se gasta en romper enlaces, \textbf{no en subir la temperatura}.
Resultado: puedes estar metiendo calor a saco y el termómetro se queda clavado en
la temperatura de fusión hasta que todo el trocito se ha fundido. Lo mismo, aún más
exagerado, al hervir.

\begin{definicion}[Calor latente]
El \emph{calor latente} de fusión $L_f$ (y el de vaporización $L_v$) es la energía
por unidad de masa que absorbe el material al cambiar de fase \textbf{a temperatura
casi constante}. Para el acero $L_f\approx\SI{2.7e5}{J/kg}$ y
$L_v\approx\SI{6.3e6}{J/kg}$ (¡vaporizar cuesta más de 20 veces que fundir!).
\end{definicion}

\subsection{El truco del $c_p$ aparente}

¿Cómo metemos ``energía que entra sin que suba $T$'' en una fórmula como
\eqref{eq:Eint}, que dice ``energía $=$ integral de $c_p$''? Con un truco muy
elegante: fingimos que, \textbf{justo en la temperatura de fusión}, el calor
específico se dispara hacia un valor enorme. Así, para pasar de un pelín por debajo
a un pelín por encima de $T_m$ hay que meter una barbaridad de energía
(precisamente $L_f$), y eso es exactamente lo que pasa de verdad.

Matemáticamente, ese ``pico'' se modela con una campana de Gauss estrecha centrada
en la temperatura de fusión $T_m$, y otra en la de ebullición $T_b$. El calor
específico \textbf{aparente} es
\begin{equation}
  \cpapp(T) \;=\; \underbrace{\cpsens}_{\text{normal}}
  \;+\; L_f\,\frac{\dd f_\ell}{\dd T}
  \;+\; L_v\,\frac{\dd f_v}{\dd T},
  \label{eq:cpapp}
\end{equation}
donde $\cpsens$ es el calor específico ``sensible'' (el normal, $\approx 600$),
y las dos campanas son
\begin{equation}
  \frac{\dd f_\ell}{\dd T}=\frac{1}{s_\ell\sqrt{\pi}}\,
    e^{-\left((T-T_m)/s_\ell\right)^2},
  \qquad
  \frac{\dd f_v}{\dd T}=\frac{1}{s_v\sqrt{\pi}}\,
    e^{-\left((T-T_b)/s_v\right)^2}.
  \label{eq:bumps}
\end{equation}
$s_\ell$ y $s_v$ son los anchos de las campanas (lo ``repartido'' que está el
cambio de fase en temperatura). La gracia es que \textbf{el área de cada campana
vale $1$}, de modo que al integrar $\cpapp$ por la zona de fusión recuperamos
exactamente $L_f$ (y por la de ebullición, $L_v$). El término de ebullición actúa
además como un \emph{termostato físico}: la energía que si no dispararía la
temperatura a valores absurdos ($\sim\!\SI{e4}{K}$) se la ``traga'' la campana, y
la temperatura se queda contenida cerca del punto de ebullición.

\dummies{$\cpapp$ es el $c_p$ normal \emph{más} dos chepas: una gigante en la
temperatura de fusión y otra en la de ebullición. Esas chepas representan la
energía que el cambio de fase se traga sin subir la temperatura. El área de cada
chepa es justo el calor latente.}

\subsection{La entalpía: la energía almacenada, de verdad}
\label{sec:entalpia}

Ahora juntamos todo. La energía por unidad de volumen almacenada en el metal,
\emph{contando ya el calor latente}, al ir de una temperatura de referencia $T_0$
hasta $T$, es la integral del $c_p$ aparente. A esa cantidad la llamamos
\textbf{entalpía volumétrica}:
\begin{equation}
  \boxed{\;h(T) \;=\; \rho \int_{T_0}^{T} \cpapp(\tau)\,\dd\tau\;}
  \qquad [\,\mathrm{J/m^3}\,].
  \label{eq:enthalpy}
\end{equation}

\begin{idea}[El papel central de la entalpía]
$h(T)$ es la \textbf{cantidad de energía guardada por unidad de volumen}. Es la
magnitud que \emph{de verdad} se conserva y se transporta. La temperatura $T$ es
solo una \emph{etiqueta} que leemos a partir de la energía.
\end{idea}

¿Qué forma tiene la curva $h(T)$? (Figura~\ref{fig:hcurve}.)
\begin{itemize}[leftmargin=1.5em]
  \item Lejos de los cambios de fase, $\cpapp\approx\cpsens$ es casi constante, así
  que $h(T)\approx \rho\,\cpsens\,(T-T_0)$: una recta. Más energía $\to$ más
  temperatura, proporcional.
  \item En la fusión y la ebullición, $\cpapp$ se dispara (las campanas), así que
  la curva $h(T)$ se vuelve casi \textbf{vertical}: metes mucha energía ($h$ sube
  mucho) y la temperatura apenas se mueve. Son las \emph{mesetas} de cambio de
  fase.
\end{itemize}

\begin{figure}[t]
\centering
\fbox{\parbox{0.9\linewidth}{\centering\small
\textit{Curva entalpía--temperatura $h(T)$ (esquema).}\\[6pt]
\ttfamily
\begin{tabular}{l}
 $h$ (energía)\\[2pt]
 \textbar\hfill \_\_\_\_/ \ \ $\leftarrow$ meseta de ebullición ($L_v$)\\
 \textbar\hfill /\\
 \textbar\hfill \_\_\_\_/ \ \ $\leftarrow$ meseta de fusión ($L_f$, a $T\!\approx\!T_m$)\\
 \textbar\ \ \ \ \ \_\_\_\_\_\_/\\
 \textbar\ \_/ \ \ (pendiente $\rho c_p$ fuera del cambio de fase)\\
 +\rule{5cm}{0.4pt}$>$\ \ $T$ (temperatura)\\
\end{tabular}}}
\caption{La relación entalpía--temperatura. En las mesetas, un gran aumento de
energía $h$ produce un aumento casi nulo de temperatura $T$: el cambio de fase se
``come'' la energía. Esta no-linealidad es la protagonista de todo el documento.}
\label{fig:hcurve}
\end{figure}

Como $\cpapp>0$ siempre, $h(T)$ es \textbf{estrictamente creciente}: a cada
temperatura le corresponde una entalpía y viceversa. Por eso la podemos
\textbf{invertir}: dada la energía almacenada $h$, recuperamos la temperatura con
\begin{equation}
  \boxed{\;T = h^{-1}(h).\;}
  \label{eq:hinv}
\end{equation}
Esto es lo que en física se llama el \emph{método de la entalpía}: lleva la cuenta
en energía, y traduce a temperatura solo al final.

\dummies{La entalpía $h$ es ``cuánta energía hay guardada aquí''. Su gráfica contra
la temperatura es una rampa con dos rellanos planos (fundir y hervir). Como sube
siempre, puedes ir de $h$ a $T$ o de $T$ a $h$ sin ambigüedad: $T=h^{-1}(h)$.}

%==============================================================================
\section{La ecuación del calor, deducida de un balance de energía}
\label{sec:calor}
%==============================================================================

Hasta ahora hablamos de un trocito de metal. Ahora queremos una ecuación para
\emph{todo} el campo $T(\bm x,t)$. La obtenemos aplicando el primer principio a un
volumen diminuto y dejándolo encogerse.

\subsection{Dos ingredientes}

\paragraph{(a) ¿Cómo viaja el calor? Ley de Fourier.} La experiencia dice que el
calor fluye de lo caliente a lo frío, y más deprisa cuanto más brusco es el cambio
de temperatura. Eso se escribe con el \emph{flujo de calor} $\bm q$ (energía que
cruza una superficie por unidad de área y tiempo, $\mathrm{W/m^2}$):
\begin{equation}
  \bm q = -\,k\,\grad T.
  \label{eq:fourier}
\end{equation}
Aquí $k$ es la \textbf{conductividad térmica} ($\mathrm{W\,m^{-1}\,K^{-1}}$; para
el acero $k\approx 30$) y $\grad T$ es el \textbf{gradiente} de temperatura: un
vector que apunta hacia donde $T$ crece más rápido. El signo menos dice ``el calor
va en sentido contrario al gradiente'', es decir, \textbf{hacia donde hace más
frío}. Esto es el segundo principio en versión local.

\paragraph{(b) Balance de energía.} Tomemos una cajita de volumen $\dd V$. El
primer principio dice: \emph{el ritmo al que aumenta su energía $=$ lo que entra
neto por las paredes $+$ lo que genera la fuente interna}. ``Lo que entra neto por
las paredes'' es menos la divergencia del flujo, $-\grad\!\cdot\!\bm q$ (la
divergencia mide cuánto ``sale''); y la fuente interna (el arco) la llamamos $Q$
(potencia por unidad de volumen, $\mathrm{W/m^3}$). La energía por volumen es la
entalpía $h$, cuyo ritmo de cambio es $\partial h/\partial t$. Entonces:
\begin{equation}
  \frac{\partial h}{\partial t} = -\,\grad\!\cdot\!\bm q \;+\; Q.
  \label{eq:balance}
\end{equation}

\subsection{Juntando los ingredientes}

Sustituimos Fourier \eqref{eq:fourier} en el balance \eqref{eq:balance}:
\begin{equation}
  \frac{\partial h}{\partial t} = \grad\!\cdot\!\big(k\,\grad T\big) + Q.
  \label{eq:heath}
\end{equation}
Y ahora usamos la entalpía \eqref{eq:enthalpy}. Por la regla de la cadena,
$\partial h/\partial t = \dfrac{\dd h}{\dd T}\dfrac{\partial T}{\partial t}
= \rho\,\cpapp(T)\,\dfrac{\partial T}{\partial t}$ (recuerda
$\dd h/\dd T = \rho\,\cpapp$). Sustituyendo obtenemos la forma clásica de la
\textbf{ecuación del calor}:
\begin{equation}
  \boxed{\;\rho\,\cpapp(T)\,\frac{\partial T}{\partial t}
    = \grad\!\cdot\!\big(k(T)\,\grad T\big) + Q(\bm x,t).\;}
  \label{eq:heat}
\end{equation}

\dummies{Esta ecuación solo dice ``el primer principio en cada punto'': lo que se
calienta un trocito ($\rho\cpapp\,\partial T/\partial t$) es igual a lo que le llega
por conducción de los vecinos ($\grad\!\cdot\!(k\grad T)$) más lo que le mete el
arco ($Q$). Fíjate que escrita en $T$ aparece el factor $\rho\cpapp$ delante; esto
será importantísimo más adelante.}

\subsection{La fuente del arco (Goldak) y las fronteras}

La fuente $Q$ del arco se modela con una distribución que se usa estándar en
soldadura, la \textbf{doble elipsoide de Goldak}: una ``mancha'' de calor con forma
de dos medios-balones (uno delante y otro detrás del punto de la antorcha) que se
mueve con ella. No hace falta su fórmula exacta para entender el resto; basta saber
que $Q(\bm x,t)$ es \textbf{grande y concentrada cerca de la antorcha} y cero lejos
(y cero del todo cuando el arco se apaga en la fase de enfriamiento).

En los \textbf{bordes} de la placa hay que decir qué pasa con el calor. Tres
opciones (condiciones de contorno):
\begin{itemize}[leftmargin=1.5em]
  \item \textbf{Dirichlet:} el borde se mantiene a una temperatura fija $T=T_D$
  (como si estuviera pegado a algo que no se calienta).
  \item \textbf{Convección:} el borde cede calor al aire a su alrededor, a un ritmo
  $h_{\mathrm{conv}}(T-\Tinf)$, donde $\Tinf$ es la temperatura ambiente y
  $h_{\mathrm{conv}}$ un coeficiente de intercambio.
  \item \textbf{Radiación:} el borde irradia calor, a un ritmo
  $\varepsilon\,\sigma_{\mathrm{SB}}(T^4-\Tinf^4)$, con $\varepsilon$ la emisividad
  y $\sigma_{\mathrm{SB}}$ la constante de Stefan--Boltzmann.
\end{itemize}
Convección $+$ radiación juntas se llaman condición de \emph{Robin}. No nos
detenemos más: lo que importa es que \textbf{hay un baño térmico} (el ambiente) al
que la placa tiende, y que la temperatura nunca debería bajar de ese ambiente solo
enfriándose (esto será una prueba de cordura más adelante).

%==============================================================================
\section{El método de los elementos finitos (FEM), sin misterio}
\label{sec:fem}
%==============================================================================

\subsection{El problema y la idea}

La ecuación \eqref{eq:heat} es una ecuación en derivadas parciales: nos pide $T$ en
\emph{infinitos} puntos de una placa de forma cualquiera. No sabemos resolverla con
lápiz y papel salvo en casos triviales. El FEM es la receta para resolverla
\textbf{de forma aproximada en un ordenador}.

\begin{idea}[La idea del FEM en una frase]
En vez de buscar la temperatura en infinitos puntos, partimos la placa en muchos
\textbf{triangulitos} (los \emph{elementos}), guardamos la temperatura solo en sus
\textbf{vértices} (los \emph{nodos}), y suponemos que dentro de cada triángulo la
temperatura varía de forma sencilla (lineal). Así, el campo continuo se reduce a
una \textbf{lista finita de números} $\bm T = (T_1,\dots,T_N)$: la temperatura en
cada nodo.
\end{idea}

A esa partición en triángulos se la llama \textbf{malla}. Aproximamos
\begin{equation}
  T(\bm x,t) \approx \sum_{i=1}^{N} T_i(t)\,\phi_i(\bm x),
  \label{eq:interp}
\end{equation}
donde $\phi_i(\bm x)$ son las \emph{funciones de forma}: $\phi_i$ vale $1$ en el
nodo $i$, $0$ en los demás nodos, y varía linealmente entre medias. Cada $T_i(t)$ es
la temperatura del nodo $i$ en el instante $t$, lo que queremos calcular.

\subsection{De la PDE a un sistema de ecuaciones}

Si metes la aproximación \eqref{eq:interp} en la ecuación del calor
\eqref{eq:heat} y haces el procedimiento estándar del FEM (multiplicar por una
función de prueba e integrar sobre la placa; los detalles no son necesarios aquí),
la ecuación en derivadas parciales se convierte en un \textbf{sistema de ecuaciones
diferenciales ordinarias} para la lista de temperaturas $\bm T(t)$:
\begin{equation}
  \bm M(\bm T)\,\dot{\bm T} \;+\; \bm A(\bm T)\,\bm T \;=\; \bm b(t).
  \label{eq:semidisc}
\end{equation}
Vamos pieza a pieza:
\begin{itemize}[leftmargin=1.5em]
  \item $\dot{\bm T}=\dd\bm T/\dd t$ es la velocidad de cambio de las temperaturas
  nodales.
  \item $\bm M$ es la \textbf{matriz de masa} (o de capacidad térmica). Recoge el
  término $\rho\,\cpapp$: cuánto cuesta calentar. Depende de $\bm T$ porque
  $\cpapp$ depende de $T$ (¡el calor latente!).
  \item $\bm A=\bm K+\bm R$ es la \textbf{matriz de rigidez}: $\bm K$ viene de la
  conducción ($k\grad T$) y $\bm R$ de las fronteras (convección/radiación).
  \item $\bm b(t)$ es el \textbf{vector de cargas}: la energía que mete el arco
  (Goldak) más lo que entra/sale por los bordes.
\end{itemize}

\dummies{El FEM convierte ``resuelve esta ecuación en toda la placa'' en
``resuelve este sistema de ecuaciones con una temperatura por nodo''. Las matrices
$\bm M$ (calentar cuesta) y $\bm A$ (el calor se reparte y se escapa por los
bordes) son solo la contabilidad de la energía nodo a nodo.}

\subsection{Avanzar en el tiempo}

\eqref{eq:semidisc} dice cómo cambian las temperaturas, pero queremos pasar del
instante $n$ al $n+1$. Discretizamos el tiempo en pasos $\Delta t$ y usamos un
esquema implícito (Euler hacia atrás), que es estable:
\begin{equation}
  \Big(\tfrac{1}{\Delta t}\bm M + \bm A\Big)\bm T^{\,n+1}
  = \tfrac{1}{\Delta t}\bm M\,\bm T^{\,n} + \bm b^{\,n+1}.
  \label{eq:beuler}
\end{equation}
Esto es un sistema lineal $\bm{\mathcal A}\,\bm T^{n+1}=\bm c$ que el ordenador
resuelve en cada paso. Pero hay una pega: $\bm M$ y $\bm A$ \emph{dependen} de la
$\bm T^{n+1}$ que aún no conocemos (porque $\cpapp(T)$ y la radiación $T^4$ son
no-lineales). Se arregla \textbf{iterando}: adivino $\bm T^{n+1}$, recalculo
$\bm M,\bm A$, resuelvo, y repito hasta que deja de cambiar. Eso es la
\emph{iteración de Picard}.

El resultado del FEM es una \textbf{película}: una secuencia de fotogramas
$\bm T^0,\bm T^1,\bm T^2,\dots$ con la temperatura de cada nodo en cada instante.
Esta es nuestra ``verdad de referencia'' (\emph{ground truth}) cara de calcular:
puede tardar minutos u horas. El objetivo del proyecto es \textbf{aprender un
sustituto rápido} (una red neuronal sobre la malla) que imite esta película. Cómo
se aprende ese sustituto es otra historia; aquí nos centramos en \textbf{qué física
y qué garantías} le imponemos.

%==============================================================================
\section{GENERIC desde cero: meter los dos principios en la dinámica}
\label{sec:generic}
%==============================================================================

\subsection{La bisagra energía--entropía--temperatura}
\label{sec:bisagra}

Recuperemos la relación \eqref{eq:bisagra} y justifiquémosla. El segundo principio
introduce la entropía con $\dd S = \delta Q_{\mathrm{rev}}/T$. Si a un trocito de
material, sin que haga trabajo, le añadimos calor $\delta Q$, toda esa energía se
queda como energía interna: $\dd E = \delta Q$. Por tanto
\begin{equation}
  \dd S = \frac{\dd E}{T}
  \qquad\Longrightarrow\qquad
  \frac{\partial S}{\partial E} = \frac{1}{T}.
  \label{eq:mu}
\end{equation}
A la derivada $\partial S/\partial E$ la llamaremos $\mu$ (la ``fuerza
termodinámica''):
\begin{equation}
  \mu \;\equiv\; \frac{\partial S}{\partial E} \;=\; \frac{1}{T}.
  \label{eq:mudef}
\end{equation}

\begin{idea}[Lo único que hay que recordar de aquí]
La derivada de la entropía respecto de la energía es $1/T$. Esto vale
\textbf{siempre}, con o sin calor latente, sea cual sea $c_p$, porque sale
directamente de la definición de entropía. Es robusto como una roca.
\end{idea}

\subsection{La estructura GENERIC}

GENERIC (siglas en inglés de ``Ecuación General para el Acoplamiento
No-Equilibrio Reversible--Irreversible'') es una \textbf{plantilla} para escribir
cómo evoluciona un sistema de manera que \emph{automáticamente} cumpla los dos
principios. La idea: todo proceso tiene una parte \textbf{reversible} (sin
pérdidas, como un péndulo ideal) y una parte \textbf{irreversible} (con pérdidas,
como el rozamiento que genera calor). GENERIC las separa y las controla.

Sea $\bm z$ el \textbf{estado} del sistema (la lista de variables que lo
describen). GENERIC postula:
\begin{equation}
  \boxed{\;
  \frac{\dd \bm z}{\dd t}
  = \underbrace{\bm L\,\frac{\partial E}{\partial \bm z}}_{\text{reversible}}
  + \underbrace{\bm M\,\frac{\partial S}{\partial \bm z}}_{\text{irreversible}}
  \;}
  \label{eq:generic}
\end{equation}
con dos ``motores'': la energía total $E(\bm z)$ y la entropía total $S(\bm z)$.
Las derivadas $\partial E/\partial\bm z$ y $\partial S/\partial\bm z$ son los
gradientes (las ``pendientes'') de esos potenciales. Las matrices $\bm L$ y $\bm M$
deben cumplir:
\begin{itemize}[leftmargin=1.5em]
  \item $\bm L$ es \textbf{antisimétrica}: $\bm L^{\mathsf T}=-\bm L$. (La parte
  reversible.)
  \item $\bm M$ es \textbf{simétrica y semidefinida positiva} (SPSD):
  $\bm M^{\mathsf T}=\bm M$ y $\bm x^{\mathsf T}\bm M\bm x\ge 0$ para todo $\bm x$.
  (La parte irreversible.)
  \item \textbf{Condiciones de degeneración:}
  \begin{equation}
    \bm L\,\frac{\partial S}{\partial \bm z}=\bm 0,
    \qquad
    \bm M\,\frac{\partial E}{\partial \bm z}=\bm 0.
    \label{eq:degen}
  \end{equation}
  En palabras: la parte reversible \emph{no toca la entropía}, y la parte
  irreversible \emph{no toca la energía}.
\end{itemize}

\subsection{El milagro: los dos principios salen solos}

Veamos que \eqref{eq:generic}--\eqref{eq:degen} garantizan los principios.

\paragraph{Primer principio ($\dd E/\dd t = 0$ si nadie aporta energía).} La
energía cambia según la regla de la cadena, $\dd E/\dd t = (\partial
E/\partial\bm z)^{\mathsf T}\,\dd\bm z/\dd t$. Sustituyendo \eqref{eq:generic}:
\begin{equation}
  \frac{\dd E}{\dd t}
  = \Big(\tfrac{\partial E}{\partial\bm z}\Big)^{\!\mathsf T}\bm L\,
    \tfrac{\partial E}{\partial\bm z}
  + \Big(\tfrac{\partial E}{\partial\bm z}\Big)^{\!\mathsf T}\bm M\,
    \tfrac{\partial S}{\partial\bm z}.
\end{equation}
El primer término es \textbf{cero} porque $\bm L$ es antisimétrica (para cualquier
vector $\bm a$, $\bm a^{\mathsf T}\bm L\bm a=0$: un número igual a su propio
negativo solo puede ser $0$). El segundo término es \textbf{cero} por la
degeneración $\bm M\,\partial E/\partial\bm z=\bm 0$. Total: $\dd E/\dd t=0$.

\paragraph{Segundo principio ($\dd S/\dd t \ge 0$).} Igual con la entropía:
\begin{equation}
  \frac{\dd S}{\dd t}
  = \Big(\tfrac{\partial S}{\partial\bm z}\Big)^{\!\mathsf T}\bm L\,
    \tfrac{\partial E}{\partial\bm z}
  + \Big(\tfrac{\partial S}{\partial\bm z}\Big)^{\!\mathsf T}\bm M\,
    \tfrac{\partial S}{\partial\bm z}.
\end{equation}
El primer término es \textbf{cero} por la degeneración $\bm L\,\partial
S/\partial\bm z=\bm 0$. El segundo es \textbf{$\ge 0$} porque $\bm M$ es
semidefinida positiva (eso significa exactamente que $\bm a^{\mathsf T}\bm M\bm
a\ge0$). Total: $\dd S/\dd t\ge0$.

\begin{idea}[Por qué nos encanta GENERIC]
Si construyes la dinámica con esta plantilla, \textbf{es imposible} que viole el
primer o el segundo principio, \emph{pase lo que pase} con los detalles. La energía
se conserva y la entropía no decrece \emph{por construcción}, no por suerte ni por
entrenar bien. Para una red neuronal que da millones de pasos, esto es oro: las
garantías no dependen de cuántos pasos des.
\end{idea}

\subsection{Nuestro caso: conducción pura y sistema abierto}

La conducción de calor \textbf{no tiene parte reversible} (no hay un ``flujo de
calor que vaya y vuelva sin pérdidas''): es puro rozamiento térmico. Por tanto
$\bm L=\bm 0$ y nos quedamos solo con la parte irreversible:
\begin{equation}
  \frac{\dd\bm z}{\dd t} = \bm M\,\frac{\partial S}{\partial\bm z}.
  \label{eq:diss}
\end{equation}

Una sutileza honesta: la placa que soldamos \textbf{no} está aislada: el arco le
mete energía y los bordes la dejan escapar. La identidad $\dd E/\dd t=0$ vale para
un sistema \emph{aislado}. La solución estándar (extensión de GENERIC a sistemas
abiertos) es: la estructura SPSD gobierna solo la \textbf{evolución interna} (la
conducción entre nodos), y el intercambio con el exterior (arco $+$ bordes) se añade
como un término aparte $\bm f^{\mathrm{ext}}$ que \emph{no} está obligado a
conservar energía (¡faltaría más, si es justo el que la mete o la saca!):
\begin{equation}
  \frac{\dd\bm z}{\dd t}
  = \underbrace{\bm M\,\frac{\partial S}{\partial\bm z}}_{\text{interno: conserva energía, } \dd S\ge0}
  \;+\; \underbrace{\bm f^{\mathrm{ext}}}_{\text{arco + bordes}}.
  \label{eq:open}
\end{equation}
Así, la energía interna se conserva exactamente y la energía total cambia
\emph{justo} en lo que entra por el arco o sale por los bordes, que es lo que la
física pide.

%==============================================================================
\section{El error de usar la temperatura como estado}
\label{sec:error}
%==============================================================================

Llega el momento clave del documento. Tenemos la plantilla \eqref{eq:diss}. Falta
decidir \textbf{cuál es el estado $\bm z$}. La tentación natural es:
\textbf{``pues la temperatura, $\bm z=\bm T$''}. Vamos a ver que eso está
\emph{mal} en cuanto hay calor latente.

\subsection{Cómo conectamos los nodos: el grafo de conductancias}

Antes, necesitamos darle forma a $\bm M$ en una malla. Pensamos los nodos como un
\textbf{grafo}: cada nodo está unido a sus vecinos por aristas. A cada arista
$i\!-\!j$ le asignamos una \textbf{conductancia} $w_{ij}\ge 0$ (cómo de bien pasa el
calor por ahí), simétrica $w_{ij}=w_{ji}$. Con eso se construye la matriz
\textbf{Laplaciana} del grafo,
\begin{equation}
  \Lw = \bm D - \bm W,\qquad
  \bm W=[w_{ij}],\quad \bm D=\mathrm{diag}\Big(\textstyle\sum_j w_{ij}\Big).
  \label{eq:laplacian}
\end{equation}
Esta $\Lw$ tiene dos propiedades que vienen como anillo al dedo:
\begin{enumerate}[leftmargin=1.5em]
  \item Es \textbf{SPSD} (simétrica y $\ge 0$), justo lo que GENERIC pide para
  $\bm M$.
  \item \textbf{Suma cero por filas}: $\Lw\bm 1=\bm 0$ (donde $\bm 1$ es el vector
  de unos). Esto va a ser \emph{exactamente} la condición de degeneración
  $\bm M\,\partial E/\partial\bm z=\bm0$.
\end{enumerate}
Aplicada a un vector $\bmu$ de ``fuerzas'' nodales, la Laplaciana da, en cada nodo,
\begin{equation}
  (\Lw\bmu)_i = \sum_{j\sim i} w_{ij}\,(\mu_i-\mu_j),
  \label{eq:lapaction}
\end{equation}
es decir, la suma de las diferencias con los vecinos, ponderadas por la
conductancia. (El símbolo $j\sim i$ significa ``$j$ vecino de $i$''.)

\subsection{La tentación (y por qué falla)}

Tomemos $\bm z=\bm T$ y, copiando \eqref{eq:diss}, escribamos directamente el
incremento de temperatura como
\begin{equation}
  \Delta T_i = (\Lw\bmu)_i,\qquad \mu_i=\frac{1}{T_i}.
  \label{eq:wrong}
\end{equation}
Sumando sobre todos los nodos, como $\Lw$ suma cero por filas,
\begin{equation}
  \sum_i \Delta T_i = \bm 1^{\mathsf T}\Lw\bmu = 0.
  \label{eq:sumT}
\end{equation}
``¡Bien, la suma de temperaturas se conserva, esto conserva la energía!'' \alert{Mal.}
Conservar $\sum_i T_i$ \textbf{no es} conservar la energía. La energía de un nodo no
es su temperatura: es su entalpía. El cambio de energía del nodo $i$ es
$\Delta e_i = \rho\,\cpapp(T_i)\,\Delta T_i$ (recuerda $\dd h=\rho\cpapp\,\dd T$).
Por tanto la \emph{verdadera} conservación de energía sería
\begin{equation}
  \sum_i \Delta e_i = \sum_i \rho\,\cpapp(T_i)\,\Delta T_i = 0,
  \label{eq:energytrue}
\end{equation}
y esto \textbf{no} coincide con \eqref{eq:sumT} salvo que $\cpapp$ sea el mismo en
todos los nodos. Las dos sumas se diferencian justo en el factor $\cpapp(T_i)$,
\textbf{que se dispara en el baño de fusión}.

\begin{equation}
  \underbrace{\sum_i \Delta T_i = 0}_{\text{suma de temperaturas}}
  \;\;\neq\;\;
  \underbrace{\sum_i \rho\,\cpapp(T_i)\,\Delta T_i = 0}_{\text{energía (lo correcto)}}.
  \label{eq:wrongcons}
\end{equation}

\begin{ejemplo}[Dos nodos, para verlo con números]
Nodo $A$ en plena fusión, con $\cpapp$ enorme (digamos $20$ veces el normal por la
campana de $L_f$); nodo $B$ frío, con $\cpapp$ normal. La dinámica correcta mueve
una cierta energía $\Delta e$ de $A$ a $B$: $\Delta e_A=-\Delta e$,
$\Delta e_B=+\Delta e$, y \textbf{la energía se conserva} ($\Delta e_A+\Delta e_B=0$).
Traducido a temperaturas: $\Delta T_A=\Delta e_A/(\rho\,(\cpapp)_A V)$ es
\emph{minúsculo} (el $(\cpapp)_A$ enorme en el denominador), mientras que
$\Delta T_B=\Delta e_B/(\rho\,(\cpapp)_B V)$ es $20$ veces mayor. ¡La suma
$\Delta T_A+\Delta T_B\neq0$! Si en cambio \emph{forzaras} $\sum\Delta T=0$ como en
\eqref{eq:wrong}, estarías moviendo mal la energía: enfriarías $A$ de más
(le restarías mucha más temperatura de la que la física permite, porque ignoras que
ahí cada grado cuesta $20$ veces más).
\end{ejemplo}

Físicamente, \eqref{eq:energytrue} reescrita como
\begin{equation}
  \frac{\dd T_i}{\dd t} = \frac{1}{\rho\,\cpapp(T_i)}\,\frac{\dd e_i}{\dd t}
  \label{eq:correctT}
\end{equation}
dice que \textbf{en el baño de fusión $\dd T/\dd t$ debe ser pequeño} (toda la
energía que llega se la traga el cambio de fase). La versión ingenua
\eqref{eq:wrong} \emph{no tiene} ese freno $1/(\rho\cpapp)$ y por eso
\textbf{mueve de más los nodos calientes}: achata o deforma el pico de temperatura.
Esto explica lo que se observaba experimentalmente: el modelo con estructura sobre
$T$ ``capaba'' el pico de fusión.

\dummies{Conservar la \emph{suma de temperaturas} no es conservar la \emph{energía},
porque cada grado vale distinto según dónde estés en la curva $h(T)$ (en la fusión,
un grado cuesta un montón de julios). Imponer GENERIC sobre $T$ comete justo este
error, y se nota más donde más nos importa: en el charco de fusión.}

%==============================================================================
\section{La solución: GENERIC entálpico}
\label{sec:enthalpic}
%==============================================================================

El arreglo es de cajón una vez visto el problema: \textbf{que el estado sea la
energía, no la temperatura}. Es el ``método de la entalpía'' de toda la vida,
metido dentro de GENERIC.

\subsection{Estado, potenciales y los dos principios}

Tomamos como estado la \textbf{energía nodal} (extensiva, en julios):
\begin{equation}
  z_i = \mathcal E_i = h_i\,V_i,
  \label{eq:state}
\end{equation}
donde $h_i=h(T_i)$ es la entalpía volumétrica del nodo \eqref{eq:enthalpy} y $V_i$
es el ``volumen tributario'' del nodo (el cachito de placa que le corresponde). La
energía total es la suma, $E=\sum_i\mathcal E_i$, así que su gradiente es
trivial:
\begin{equation}
  \frac{\partial E}{\partial\mathcal E_i}=1
  \qquad\Longrightarrow\qquad
  \frac{\partial E}{\partial\bm z}=\bm 1 .
  \label{eq:gradE}
\end{equation}
¡Mira qué bonito! La condición de degeneración $\bm M\,\partial E/\partial\bm z=\bm0$
se vuelve $\Lw\bm 1=\bm 0$, que la Laplaciana \textbf{cumple exactamente} por
construcción (suma cero por filas). La energía se conservará \emph{en cualquier
malla}.

El gradiente de la entropía respecto de la energía es nuestra vieja amiga la
bisagra \eqref{eq:mudef}: $\mu_i=\partial S/\partial\mathcal E_i=1/T_i$. Con
$\bm M=\Lw$, la evolución \eqref{eq:diss} se concreta en:
\begin{equation}
  \boxed{\;
  \frac{\dd\mathcal E_i}{\dd t} = (\Lw\bmu)_i
  = \sum_{j\sim i} w_{ij}\Big(\frac{1}{T_i}-\frac{1}{T_j}\Big),
  \qquad
  T_i = h^{-1}\!\big(\mathcal E_i/V_i\big).
  \;}
  \label{eq:enthevol}
\end{equation}

Leamos esta ecuación despacio, que es \emph{la} ecuación del documento:
\begin{itemize}[leftmargin=1.5em]
  \item La parte izquierda es ``cómo cambia la energía del nodo $i$''.
  \item La derecha suma, sobre los vecinos, las diferencias de $\mu=1/T$. Si el
  nodo $i$ está \textbf{más caliente} que su vecino ($T_i>T_j\Rightarrow 1/T_i <
  1/T_j$), el término $(1/T_i-1/T_j)$ es \textbf{negativo}: el nodo caliente
  \textbf{pierde} energía. ¡El calor va de caliente a frío \emph{por
  construcción}, segundo principio en acción!
  \item Lo que pierde $i$ lo gana $j$ exactamente (la arista aporta $+$ a uno y $-$
  al otro). Energía conservada.
  \item Al final, traducimos energía a temperatura con $T_i=h^{-1}(\mathcal
  E_i/V_i)$. \textbf{Aquí entra gratis el calor latente:} en la meseta de fusión
  $h^{-1}$ es casi plano, así que un gran $\Delta\mathcal E$ produce un
  $\Delta T$ minúsculo, \emph{igual que en el FEM}. El freno $1/(\rho\cpapp)$ que le
  faltaba a la versión sobre $T$ \textbf{aparece solo} al invertir la curva.
\end{itemize}

\begin{proposicion}[Los dos principios, en cualquier malla]
\label{prop:laws}
Para cualesquiera conductancias simétricas no negativas $w_{ij}=w_{ji}\ge0$, la
evolución \eqref{eq:enthevol} cumple, en cada grafo,
\begin{equation}
  \frac{\dd E}{\dd t}=\sum_i\frac{\dd\mathcal E_i}{\dd t}
  = \bm 1^{\mathsf T}\Lw\bmu = 0
  \quad(\text{1.\textsuperscript{er} principio}),
\end{equation}
\begin{equation}
  \frac{\dd S}{\dd t}=\bmu^{\mathsf T}\Lw\bmu
  = \tfrac12\sum_{i,j} w_{ij}(\mu_i-\mu_j)^2 \ge 0
  \quad(\text{2.\textsuperscript{o} principio}).
\end{equation}
\end{proposicion}
\begin{proof}
La conservación de energía sale de $\bm 1^{\mathsf T}\Lw=\bm 0$ (suma cero por
columnas). Para la entropía, la identidad
$\bmu^{\mathsf T}\Lw\bmu=\tfrac12\sum_{i,j}w_{ij}(\mu_i-\mu_j)^2$ es la forma
cuadrática estándar de la Laplaciana, y es $\ge0$ porque cada $w_{ij}\ge0$ y cada
término va al cuadrado. (Esto es justo decir que $\Lw$ es SPSD.)
\end{proof}

\dummies{Cambiando el estado de ``temperatura'' a ``energía'', la \emph{misma}
plantilla GENERIC ahora sí conserva la energía de verdad (no la suma de
temperaturas) y produce entropía $\ge0$, y encima trata el calor latente exacto al
invertir la curva $h(T)$. Es el método de la entalpía clásico, escrito como una
Laplaciana de grafo.}

\subsection{Comprobación: ¿esto es de verdad la conducción de Fourier?}

Conviene verificar que \eqref{eq:enthevol} no es un invento, sino la conducción de
siempre disfrazada. Usando $\grad T=-T^2\,\grad(1/T)=-T^2\grad\mu$, el flujo de
Fourier \eqref{eq:fourier} se escribe $\bm q=-k\grad T=k\,T^2\,\grad\mu$.
Discretizado en el grafo, $\grad\!\cdot\!\bm q$ da \emph{exactamente} la forma
$\Lw\bmu$ con conductancias $w_{ij}\sim k\,T_iT_j/\ell_{ij}^2$ (siendo $\ell_{ij}$
la distancia entre nodos). O sea: la ecuación \eqref{eq:enthevol} \textbf{es} la
ley de Fourier; lo único que la red neuronal tiene que aprender son esas
conductancias $w_{ij}$ (que dependen de la temperatura y de la geometría).

%==============================================================================
\section{Dos detalles que hacen que funcione de verdad}
\label{sec:detalles}
%==============================================================================

\subsection{Acondicionamiento numérico: la entalpía en ``unidades de kelvin''}

Un problema práctico: la entalpía volumétrica es un número \emph{enorme},
$h\sim\SI{e9}{J/m^3}$. Su inversa tiene pendiente $\partial T/\partial h =
1/(\rho\cpapp)\sim\SI{2e-7}{}$, un número diminuto. Cuando entrenas la red, ese
factor minúsculo multiplica todos los gradientes que vuelven hacia las
conductancias, y los hunde por debajo del umbral numérico del optimizador
($\varepsilon\sim\SI{e-8}{}$ en Adam): el entrenamiento \textbf{se atasca} (la
pérdida se congela y la difusión nunca aprende).

La solución es trabajar con una entalpía \textbf{re-escalada a unidades de
temperatura} (kelvin), dividiendo por $\rho\cpsens$:
\begin{equation}
  H(T) \;=\; \frac{h(T)}{\rho\,\cpsens}
  \;=\; \int_{T_0}^{T}\frac{\cpapp(\tau)}{\cpsens}\,\dd\tau .
  \label{eq:Heq}
\end{equation}
Fuera del cambio de fase $H\approx T-T_0$ (pendiente $1$, ``casi temperatura'');
en las mesetas la pendiente se dispara igual que antes. Como esto es solo
\textbf{multiplicar por una constante}, la conservación de energía no cambia
($\sum_i\Delta H_i=0\Leftrightarrow\sum_i\Delta h_i=0$), pero ahora los incrementos
son de orden ``grados'' y los gradientes de orden $1$. Con este único cambio, los
gradientes pasaron de $\sim\!\SI{e-13}{}$ a $\sim\!\SI{e-4}{}$ y el entrenamiento se
desatascó desde la primera época.

\dummies{La entalpía en julios es un numerón que rompe el entrenamiento. La
dividimos por $\rho\cpsens$ para medirla en ``grados equivalentes''. Es solo un
cambio de unidades: no toca la física, pero arregla los números.}

\subsection{La fuente enriquecida: por qué la disipación sola no basta}

Hay un detalle físico importante. La parte disipativa \eqref{eq:enthevol}
\textbf{solo puede enfriar} el nodo más caliente (el calor va de caliente a frío,
siempre). Entonces, ¿quién \emph{calienta} el pico del baño de fusión? Solo la
\textbf{fuente externa} (el arco). Con una única ganancia escalar, esa fuente era
demasiado rígida para inyectar suficiente energía en el pico, y los modelos con
estructura salían ``fríos'' en el punto más caliente.

La clave es que la fuente y los bordes son el término \emph{no conservativo}
$\bm f^{\mathrm{ext}}$ de \eqref{eq:open}, que GENERIC \textbf{no obliga} a
conservar energía. Por tanto le podemos dar más flexibilidad \emph{sin romper
ninguna garantía}. Sustituimos la ganancia escalar por una modulación aprendida
nodo a nodo:
\begin{equation}
  \Delta H_i^{\mathrm{ext}}
  = \underbrace{\mathrm{softplus}\!\big(g_q+\mathrm{MLP}_q(\cdots)\big)\,q_i}
    _{\text{calentamiento del arco }\ge 0}
  + \underbrace{\mathrm{softplus}\!\big(g_c+\mathrm{MLP}_c(\cdots)\big)\,r_i\,(\Tinf-T_i)}
    _{\text{enfriamiento de Newton en bordes }r_i}.
  \label{eq:enriched}
\end{equation}
Aquí $q_i$ es el valor de la fuente Goldak en el nodo, $r_i$ marca si el nodo está
en un borde con convección, $\mathrm{MLP}$ es una pequeña red neuronal, y
$\mathrm{softplus}(\cdot)\ge0$ garantiza que el arco solo \emph{calienta} (nunca
inventa frío). Además está ``gateado'' por $q_i$: donde no hay arco ($q_i=0$), no
puede meter energía de la nada. Las redes salen inicializadas a cero, así que el
modelo \textbf{empieza igual} que la versión con ganancia escalar y luego aprende
dónde inyectar más.

Lo esencial: \textbf{la disipación \eqref{eq:enthevol} sigue intacta y
estructura-preservante}; solo el término externo, que de todas formas no estaba
sujeto a las garantías, se vuelve flexible. Conservación de energía interna y
$\dd S/\dd t\ge0$ se mantienen palabra por palabra.

\dummies{La parte ``con garantías'' solo sabe enfriar. Para recalentar el charco
hace falta el arco. Como el arco es energía que entra de fuera (no está obligado a
conservar nada), le dejamos ser flexible sin estropear los principios.}

%==============================================================================
\section{Recapitulación: el camino completo en una página}
\label{sec:resumen}
%==============================================================================

\begin{enumerate}[leftmargin=1.6em]
  \item \textbf{Física.} Soldar mueve una fuente de calor brutal por una placa; el
  metal funde y se enfría. Queremos el campo $T(\bm x,t)$.
  \item \textbf{Energía.} Calentar cuesta $c_p$ por grado. El cambio de fase se
  traga calor latente sin subir $T$ $\Rightarrow$ lo modelamos con un $\cpapp$ con
  dos campanas \eqref{eq:cpapp}.
  \item \textbf{Entalpía.} La energía almacenada es $h(T)=\rho\int\cpapp\,\dd T$
  \eqref{eq:enthalpy}: una rampa con dos rellanos. Es invertible: $T=h^{-1}(h)$.
  Es la magnitud que de verdad se conserva.
  \item \textbf{Ecuación del calor.} Balance de energía $+$ Fourier $\Rightarrow$
  $\rho\cpapp\,\partial_t T=\grad\!\cdot\!(k\grad T)+Q$ \eqref{eq:heat}.
  \item \textbf{FEM.} Partimos la placa en triángulos, una temperatura por nodo, y
  la PDE se vuelve un sistema $\bm M\dot{\bm T}+\bm A\bm T=\bm b$ que avanzamos en
  el tiempo. Es la verdad de referencia cara.
  \item \textbf{GENERIC.} Plantilla $\dot{\bm z}=\bm L\,\partial_z E+\bm
  M\,\partial_z S$ con $\bm L$ antisimétrica, $\bm M$ SPSD y degeneraciones
  $\Rightarrow$ primer y segundo principio \emph{por construcción}
  \eqref{eq:generic}--\eqref{eq:degen}.
  \item \textbf{El error.} Con estado $=$ temperatura, $\sum\Delta T=0$ \emph{no es}
  conservar energía cuando hay calor latente \eqref{eq:wrongcons}; achata el pico.
  \item \textbf{La cura.} Estado $=$ energía $\mathcal E_i=h_iV_i$. La misma
  Laplaciana de conductancias conserva energía de verdad, produce entropía, y trata
  el latente exacto al hacer $T=h^{-1}$ \eqref{eq:enthevol}, Prop.~\ref{prop:laws}.
  \item \textbf{Detalles.} Re-escalar a $H=h/(\rho\cpsens)$ para que entrene
  \eqref{eq:Heq}; fuente externa flexible \eqref{eq:enriched} para recalentar el
  pico sin romper las garantías.
\end{enumerate}

\begin{idea}[La moraleja, en una frase]
Impón la estructura termodinámica sobre la \textbf{energía} (entalpía), nunca sobre
la temperatura: bajo calor latente, conservar la suma de temperaturas no es
conservar la energía, y el desajuste es máximo justo en el baño de fusión, que es
lo que más nos importa.
\end{idea}

%==============================================================================
\appendix
\section{Glosario de símbolos}
%==============================================================================

\renewcommand{\arraystretch}{1.25}
\noindent
\begin{tabular}{@{}p{0.20\linewidth}p{0.74\linewidth}@{}}
\toprule
Símbolo & Significado \\
\midrule
$T,\ \Tinf$        & temperatura (K); temperatura ambiente \\
$\rho$             & densidad ($\mathrm{kg/m^3}$); acero $7850$ \\
$\cpsens$          & calor específico ``normal'' (sensible), $\approx 600\ \mathrm{J\,kg^{-1}K^{-1}}$ \\
$\cpapp(T)$        & calor específico aparente: $\cpsens$ + campanas de calor latente \eqref{eq:cpapp} \\
$L_f,\ L_v$        & calor latente de fusión y de vaporización (J/kg) \\
$T_m,\ T_b$        & temperatura de fusión y de ebullición (K) \\
$k$                & conductividad térmica ($\mathrm{W\,m^{-1}K^{-1}}$); acero $\approx 30$ \\
$h(T)$             & entalpía volumétrica $=\rho\int\cpapp\,\dd T$ ($\mathrm{J/m^3}$), energía almacenada \\
$H(T)$             & entalpía re-escalada a kelvin, $h/(\rho\cpsens)$ \eqref{eq:Heq} \\
$\bm q$            & flujo de calor ($\mathrm{W/m^2}$); Fourier $\bm q=-k\grad T$ \\
$Q$                & fuente volumétrica del arco ($\mathrm{W/m^3}$), Goldak \\
$\grad,\ \grad\!\cdot$ & gradiente y divergencia \\
$\bm T$            & lista de temperaturas nodales $(T_1,\dots,T_N)$ \\
$\bm M,\bm A,\bm b$ & matrices de masa, rigidez y vector de cargas del FEM \eqref{eq:semidisc} \\
$\bm z$            & estado del sistema en GENERIC \\
$E,\ S$            & energía total y entropía total (potenciales de GENERIC) \\
$\mu_i=1/T_i$      & fuerza termodinámica, $\partial S/\partial\mathcal E$ \eqref{eq:mudef} \\
$\bm L,\ \bm M$    & operador reversible (antisimétrico) e irreversible (SPSD) \\
$w_{ij}$           & conductancia de la arista $i\!-\!j$ ($\ge0$, simétrica) \\
$\Lw=\bm D-\bm W$  & Laplaciana de grafo (hace de $\bm M$); SPSD y $\Lw\bm1=\bm0$ \\
$\mathcal E_i=h_iV_i$ & energía nodal (estado en la versión entálpica) \\
$V_i$              & volumen tributario del nodo $i$ \\
$\bm f^{\mathrm{ext}}$ & término externo no conservativo (arco + bordes) \\
\bottomrule
\end{tabular}

\end{document}
